\subsection{REGULAR ELEMENTS OF A LIE ALGEBRA}

{[}Recall that \(k\) is assumed to be infinite from now on.{]}

Let g be a Lie algebra of dimension n.~If \(x \in g\) , write the characteristic polynomial of ad x in the form

\[
\det (\mathrm{T} - \operatorname{ad} x) = \sum_ {i = 0} ^ {n} a _ {i} (x) \mathrm{T} ^ {i}, \quad \text { with } a _ {i} (x) \in k.
\]

We have \(a_i(x) = (-1)^{n - i}\mathrm{Tr}\left(\bigwedge^{n - i}\operatorname {ad}x\right)\) , cf.~Algebra, Chap. III, §8, no. 11. This shows that \(x\mapsto a_i(x)\) is a homogeneous polynomial map of degree \(n - i\) from \(\mathfrak{g}\) to \(k\) (Algebra, Chap. IV, §5, no. 9).

Remarks. 1) If \(\mathfrak{g} \neq \{0\}\) , \(a_0 = 0\) since (ad \(x(x) = 0\) for all \(x \in \mathfrak{g}\) .

\begin{enumerate}
\def\labelenumi{\arabic{enumi})}
\setcounter{enumi}{1}
\tightlist
\item
  Let \(k'\) be an extension of \(k\) . Write \(\det(\mathrm{T} - \operatorname{ad} x') = \sum_{i=0}^{n} a_i'(x') \mathrm{T}^i\) for \(x' \in \mathfrak{g} \otimes_k k'\) . Then \(a_i'| \mathfrak{g} = a_i\) for all \(i\) .
\end{enumerate}

DEFINITION 2. The rank of g, denoted by rk(g), is the smallest integer l such that \(a_{l} \neq 0\) . An element x of g is called regular if \(a_{l}(x) \neq 0\) .

For all \(x \in \mathfrak{g}\) , \(\operatorname{rk}(\mathfrak{g}) \leq \dim \mathfrak{g}^0(x)\) , and equality holds if and only if \(x\) is regular.

The set of regular elements is dense and open in \(\mathfrak{g}\) for the Zariski topology (App. I).

Examples. 1) If g is nilpotent, rk(g) = dim g and all elements of g are regular.

\begin{enumerate}
\def\labelenumi{\arabic{enumi})}
\setcounter{enumi}{1}
\tightlist
\item
  Let \(\mathfrak{g} = \mathfrak{sl}(2,k)\) . If \(x = \begin{pmatrix} \gamma & \alpha \\ \beta & -\gamma \end{pmatrix} \in \mathfrak{g}\) , an easy calculation gives
\end{enumerate}

\[
\det (\mathrm{T} - \operatorname{ad} x) = \mathrm{T} ^ {3} - 4 (\alpha \beta + \gamma^ {2}) \mathrm{T}.
\]

If the characteristic of \(k\) is \(\neq 2\) , then \(\mathrm{rk}(\mathfrak{g}) = 1\) and the regular elements are those \(x\) such that \(\alpha \beta + \gamma^2 \neq 0\) .

\begin{enumerate}
\def\labelenumi{\arabic{enumi})}
\setcounter{enumi}{2}
\tightlist
\item
  Let V be a vector space of finite dimension n, and \(\mathfrak{g} = \mathfrak{gl}(V)\) . Let \(\overline{k}\) be an algebraic closure of k. Let \(x \in g\) , and let \(\lambda_{1}, \ldots, \lambda_{n}\) be the roots in \(\overline{k}\) of the characteristic polynomial of x (each root being written a number of times equal to its multiplicity). The canonical isomorphism from \(V^{*} \otimes V\) to g is compatible with the g-module structures of these two spaces, in other words it takes \(1 \otimes x - {}^{t}x \otimes 1\) to ad x (Chap. I, §3, no. 3, Prop. 4). In view of §1, Prop. 4 (i), it follows that the roots of the characteristic polynomial of ad x are the \(\lambda_{i} - \lambda_{j}\) for \(1 \leq i \leq n, 1 \leq j \leq n\) (each root being written a number of times equal to its multiplicity). Thus, the rank of g is n, and x is regular if and only if each \(\lambda_{i}\) is a simple root of the characteristic polynomial of x.
\end{enumerate}

PROPOSITION 6. Let \(\mathfrak{g}\) be a Lie algebra, \(k'\) an extension of \(k\) , and \(\mathfrak{g}' = \mathfrak{g} \otimes_k k'\) .

\begin{enumerate}
\def\labelenumi{(\roman{enumi})}
\tightlist
\item
  An element \(x\) of \(\mathfrak{g}\) is regular in \(\mathfrak{g}\) if and only if \(x \otimes 1\) is regular in \(\mathfrak{g}'\) .\\
\item
  \(\operatorname{rk}(\mathfrak{g}) = \operatorname{rk}(\mathfrak{g}')\) .
\end{enumerate}

This follows from Remark 2.

PROPOSITION 7. Let \((\mathfrak{g}_i)_{i\in \mathrm{I}}\) be a finite family of Lie algebras, and let \(\mathfrak{g} = \prod_{i\in \mathrm{I}}\mathfrak{g}_i\) .

\begin{enumerate}
\def\labelenumi{(\roman{enumi})}
\tightlist
\item
  An element \((x_{i})_{i\in \mathrm{I}}\) of \(\mathfrak{g}\) is regular in \(\mathfrak{g}\) if and only if, for all \(i\in \mathrm{I}\) , \(x_{i}\) is regular in \(\mathfrak{g}_i\) .
\end{enumerate}

\[
\text {(ii)} \operatorname{rk} (\mathfrak {g}) = \sum_ {i \in I} \operatorname{rk} (\mathfrak {g} _ {i}).
\]

Indeed, for any \(x = (x_{i})_{i \in \mathrm{I}} \in \mathfrak{g}\) , the characteristic polynomial of \(ad_{g}x\) is the product of the characteristic polynomials of the \(ad_{g_{i}}x_{i}\) .

PROPOSITION 8. Let \(f: \mathfrak{g} \to \mathfrak{g}'\) be a surjective homomorphism of Lie algebras.

\begin{enumerate}
\def\labelenumi{(\roman{enumi})}
\item
  If \(x\) is a regular element of \(\mathfrak{g}\) , \(f(x)\) is regular in \(\mathfrak{g}'\) . The converse is true if \(\operatorname{Ker} f\) is contained in the centre of \(\mathfrak{g}\) .
\item
  \(\operatorname{rk}(\mathfrak{g}) \geq \operatorname{rk}(\mathfrak{g}')\) .
\end{enumerate}

Put \(\mathrm{rk}(\mathfrak{g}) = r, \mathrm{rk}(\mathfrak{g}') = r'\) . Let \(x \in \mathfrak{g}\) . The characteristic polynomials of \(\operatorname{ad} x\) , \(\operatorname{ad} f(x)\) and \(\operatorname{ad} x|\operatorname{Ker} f\) are of the form

\[
\mathrm{P} (\mathrm{T}) = \mathrm{T} ^ {n} + a _ {n - 1} (x) \mathrm{T} ^ {n - 1} + \dots + a _ {r} (x) \mathrm{T} ^ {r},
\]

\[
\mathrm{Q} (\mathrm{T}) = \mathrm{T} ^ {n ^ {\prime}} + b _ {n ^ {\prime} - 1} (x) \mathrm{T} ^ {n ^ {\prime} - 1} + \dots + b _ {r ^ {\prime}} (x) \mathrm{T} ^ {r ^ {\prime}},
\]

\[
\mathrm{R} (\mathrm{T}) = \mathrm{T} ^ {n ^ {\prime \prime}} + c _ {n ^ {\prime \prime} - 1} (x) \mathrm{T} ^ {n ^ {\prime \prime} - 1} + \dots + c _ {r ^ {\prime \prime}} (x) \mathrm{T} ^ {r ^ {\prime \prime}},
\]

where the \(a_i, b_i, c_i\) are polynomial functions on \(\mathfrak{g}\) , with \(a_r \neq 0\) , \(b_{r'} \neq 0\) , \(c_{r''} \neq 0\) . We have \(\mathrm{P} = \mathrm{QR}\) , so \(r = r' + r''\) and \(a_r(x) = b_{r'}(x)c_{r''}(x)\) , which proves (ii) and the first assertion of (i). If \(\operatorname{Ker} f\) is contained in the centre of \(\mathfrak{g}\) , \(\mathrm{R(T)} = \mathrm{T}^{n''}\) and so \(a_r(x) = b_{r'}(x)\) , hence the second assertion of (i).

COROLLARY. Let \(\mathcal{C}_n\mathfrak{g}\) ( \(n \geq 0\) ) be a term of the ascending central series of \(\mathfrak{g}\) (Chap. I, §1, no. 6). The regular elements of \(\mathfrak{g}\) are those whose image in \(\mathfrak{g}/\mathcal{C}_n\mathfrak{g}\) is regular.

PROPOSITION 9. Let \(\mathfrak{g}\) be a Lie algebra, \(\mathfrak{g}'\) a subalgebra of \(\mathfrak{g}\) . Every element of \(\mathfrak{g}'\) regular in \(\mathfrak{g}\) is regular in \(\mathfrak{g}'\) .

For \(x \in g'\) , the restriction of \(ad_{g}x\) to \(g'\) is \(ad_{g'}x\) , and so defines an endomorphism \(u(x)\) of the vector space \(g/g'\) by passage to the quotient. Let \(d_{0}(x)\) (resp. \(d_{1}(x)\) ) be the dimension of the nilspace of \(ad_{g'}(x)\) (resp. of \(u(x)\) ), and let \(c_{0}\) (resp. \(c_{1}\) ) be the minimum of \(d_{0}(x)\) (resp. \(d_{1}(x)\) ) when x belongs to \(g'\) . There exist non-zero polynomial maps \(p_{0}, p_{1}\) from \(g'\) to k such that

\[
d _ {0} (x) = c _ {0} \Longleftrightarrow p _ {0} (x) \neq 0, \quad d _ {1} (x) = c _ {1} \Longleftrightarrow p _ {1} (x) \neq 0.
\]

Since k is infinite, the set S of \(x \in g'\) such that \(d_{0}(x) = c_{0}\) and \(d_{1}(x) = c_{1}\) is non-empty. Every element of S is regular in \(g'\) . On the other hand, S is the set of elements of \(g'\) such that the nilspace of \(ad_{g}x\) has minimum dimension, and thus contains every element of \(g'\) regular in g.

Remark. 3) Elements of \(\mathfrak{g}'\) regular in \(\mathfrak{g}\) do not necessarily exist. If at least one does exist, the set of these elements is precisely the set denoted by S in the above proof.

